\section{Study design and relationship to the pilot}
The unpublished pilot compared seven encoders on earlier Polyvore and IQON splits. Its aggregate test results were known when we chose the v2 roster, readouts, matched pairs, selection rule and statistical tests. The pilot results remain in the public repository history. We froze and hashed the v2 protocol before test evaluation (tag \texttt{v2-freeze}). Only POG supplies a confirmatory replication on previously unevaluated test data. Supplementary Table~S7 records amendments to the protocol and their consequences.

All encoders face the same outfits and completion candidates, so we compare their errors on paired questions. The full roster compares released systems, including differences in architecture, pretraining and preprocessing. Matched parents narrow the comparison to an adaptation and its training recipe. Varying the readout, labels, question construction and resampling groups tests how far the resulting rankings generalize.

\paragraph{Selection and replication}
We select models before examining their test ranks. Within each corpus and family we choose the eligible encoder with the highest native-probe validation AUC and use it in subsequent comparisons. Equal validation scores are resolved by checkpoint key. FITB, A100 and the contextual comparison can then assess consequences of a choice made using whole-outfit discrimination. Their labels have no role in that initial decision. The POG test offers an additional check because the benchmark construction rules and roster were fixed before its test evaluation. A100 tests another kind of generalization: a classifier fitted to one source is applied to questions and answer distributions supplied by another dataset.
